\subsection{Checklist}

\chk{Benchmark}:
%
We used 4G traces from Mahimahi’s \texttt{traces/} directory (AT\&T, T-Mobile, Verizon) and 5G traces from the CellReplay repository (T-Mobile and Verizon, with 3 different signal strengths each). CellReplay's 5G traces were used in both Mahimahi and CellReplay to evaluate ABR performance under the same conditions.

\chk{Models}:
%
We used six ABR models; BOLA, Gelato, Linear BBA, TTP, MPC, and Pensieve. Puffer provides these ABR models in its repository, ready-to-use.

\chk{Data set}:
%
We obtained our data set from Puffer's InfluxDB database. During each test run, Puffer logs video data such as quality, latency, and buffering, with the experiment number for each ABR model and timestamp.

\chk{Run-time environment}:
%
We used a virtual machine with Ubuntu 24.04. Our artifacts are supported in this environment, but due to version-specific requirements, we had to modify some of the source files mentioned in~\ref{s:artifact_abstract}.

\chk{Run-time state}:
%
Experimentation in Puffer is sensitive to background network traffic. However, we conducted all our tests inside an emulated network namespace on a virtual machine. In this way, our testbed is not affected by the background traffic from the host system.

\chk{Metrics and Output}:
%
We explain each metric in detail in Section~\ref{s:background}. To evaluate each model; SSIM index, buffer time, BufRatio, RTT, and delivery rate are used. These are found in our data set as a CSV file format. In Section~\ref{sec:tcp-analysis}, retransmission is also introduced from our TCP captures, saved in PCAP files.

\chk{Approximate disk space}:
%
Our virtual machine has a disk space of 100 GB. In the end, we used 47.6\% of the total disk space, which includes duplicate files, unused libraries, and repositories. It is safe to say that the minimum storage limit should be 50 GB.

\chk{Approximate preparation time}:
%
To setup the whole testbed can take up to two weeks. This includes setting up Puffer, CellReplay, and Mahimahi as specified in~\ref{s:artifact_abstract}, and testing to verify whether each component works correctly, and if not debugging.

\chk{Approximate time to complete experiments}:
%
We ran 90 tests in total, and each test run takes around 10 minutes. This means that to conduct all 90 tests, for six ABR models in all environments, it takes around 15 hours. If we were to include extracting and processing the dataset, it would take another two weeks to complete the experimentation.

%%% Local Variables:
%%% mode: latex
%%% TeX-master: "../thesis"
%%% End:
